\section{Data}
\label{sec:data}

\subsection{Source}

The Retailrocket e-commerce dataset~\cite{retailrocket2017dataset} is an
anonymised log of browsing behaviour on a real online store: roughly $2.7$M
events from about $1.4$M visitors over four and a half months, each row
carrying a \texttt{timestamp}, \texttt{visitorid}, \texttt{event}
(\texttt{view}, \texttt{addtocart} or \texttt{transaction}), \texttt{itemid},
and an optional \texttt{transactionid}. A companion \texttt{category\_tree.csv}
and an item-properties changelog (\texttt{timestamp, itemid, property, value})
supply a shallow category hierarchy and catalogue metadata such as
availability. The raw files are never committed to the repository or
downloaded automatically; \texttt{scripts/download\_data.py --kaggle} fetches
them given \texttt{KAGGLE\_USERNAME}/\texttt{KAGGLE\_KEY}.

\subsection{Synthetic fixture}

By default the repository runs against a synthetic fixture generated by the
same script with no Kaggle account. The generator is not a schema stub: it
reproduces the structural properties the model depends on.

\begin{itemize}
  \item \textbf{Bursty sessions, long idle gaps.} Each visitor is assigned a
        handful of browsing sessions spread across a configurable time span.
        Intra-session gaps are seconds to a few minutes; inter-session gaps
        are drawn independently across the whole span, so they almost always
        exceed the $30$-minute sessionization threshold
        (Section~\ref{sec:sessionization}) and the segmentation rule is
        exercised meaningfully rather than trivially.
  \item \textbf{Long-tailed item popularity.} Item weights are drawn from a
        Pareto distribution, so a small set of items accounts for a large
        share of views -- the regime in which a popularity baseline is
        genuinely competitive, not a strawman.
  \item \textbf{Latent taste clusters.} Visitors and items are each assigned
        to one of six latent clusters; a visitor's items are drawn with
        higher weight from their own cluster. Without this, no
        similarity-based method could recover structure above chance.
  \item \textbf{Realistic event funnel.} Within a session, an item can
        progress \texttt{view} $\to$ \texttt{addtocart} $\to$
        \texttt{transaction} with cluster-dependent conversion
        probabilities, matching the funnel shape of the real dataset.
\end{itemize}

\subsection{Sessionization}
\label{sec:sessionization}

Retailrocket carries no explicit session identifier, so one is constructed
directly from the inter-event gap. Given a visitor's events sorted by
\texttt{timestamp}, a new session starts whenever the gap since the previous
event exceeds a threshold $\tau$:

\begin{equation}
  \text{new\_session}(e_t) = \mathbb{1}\!\left[\, t_{e} - t_{e-1} > \tau \,\right],
  \qquad \tau = 30\text{ minutes},
  \label{eq:sessionization}
\end{equation}

applied independently per visitor (the first event of every visitor also
starts a new session). This $30$-minute inactivity rule is the convention
used throughout the session-based recommendation
literature~\cite{jannach2017knn,ludewig2018evaluation} and is what turns an
unordered event log into the unit session-kNN actually operates on.

\subsection{Why anonymous sessions, not visitors}

A durable \texttt{visitorid} would let a model build a long-run profile, the
way the companion BPR and SVD projects in this portfolio do for MovieLens and
the Anime Recommendations Database. Retailrocket does not support that
assumption: cookies churn, and the useful history for predicting the next
click is almost always confined to the session in progress. Section
\ref{sec:features} keeps a single, explicitly ablation-only visitor feature
to test this claim rather than merely assert it.
